% One-page summary of the EHMbrAIn study. Build: latexmk -pdf -outdir=build EHMbrAIn.tex
% All numbers trace to versioned artifacts in data/processed/ (see footer).
\documentclass[11pt,twocolumn]{article}
\usepackage[a4paper,margin=1.5cm,columnsep=7mm]{geometry}
\usepackage[T1]{fontenc}
\usepackage[utf8]{inputenc}
\usepackage{lmodern}
\usepackage{microtype}
\usepackage{booktabs}
\usepackage{amsmath}
\usepackage[hidelinks]{hyperref}
\usepackage{balance}
\usepackage{titlesec}
\titlespacing*{\section}{0pt}{5pt plus 1pt}{2pt}
\titleformat{\section}{\normalsize\bfseries}{}{0pt}{}
\setlength{\parskip}{2pt}
\setlength{\parindent}{0.9em}
\pagestyle{empty}

\begin{document}

\twocolumn[{%
\centering
{\LARGE\bfseries Does AI really improve jet-engine health monitoring?\par}
\vspace{0.5mm}
{\large We built a head-to-head test with ground truth. Here is where AI wins,
where it cannot --- and what the difference is worth.\par}
\vspace{1.5mm}
{\small\itshape One-page summary of \textbf{EHMbrAIn}: AI-based vs.\ traditional engine
health monitoring, a pre-registered gas-path-analysis testbed on the CFM56-7B with
synthetic ground truth\par}
\vspace{3mm}
}]

\balance
\section{The question, made falsifiable}
Does AI beat the traditional engine-health-monitoring stack --- or is that marketing? We
built the comparison no field study can run: a pyCycle thermodynamic twin of the CFM56-7B26
calibrated to certificate data, a 100-engine synthetic fleet (SynCFM56) with known
ground-truth health, and two \emph{tuned} practitioners --- the classical stack (Kalman GPA,
CUSUM, trending) and an AI suite (Mahalanobis, GRU, conformal) --- given identical data,
identical sensors, and a symmetric 50-trial tuning budget. Every hypothesis, threshold, and
test was frozen in twelve tagged pre-registrations \emph{before} the test set was touched;
failed gates were repaired physically, never by moving thresholds.

\section{Verdicts (single confirmatory pass, Holm-corrected)}
\begin{center}\small
\setlength{\tabcolsep}{3pt}
\begin{tabular}{@{}lccl@{}}
\toprule
Task & Trad. & AI & Verdict \\
\midrule
Detection recall (FA $1{=}1$) & 0.13 & \textbf{0.48} & AI, $p{=}.008$ \\
\quad median delay (cycles) & 6033 & \textbf{499} & 12$\times$ earlier \\
Isolation, confusable faults & 31\,\% & 31\,\% & \emph{tie: wall} \\
RUL RMSE @90\,\% life (cy) & 1981 & \textbf{858} & AI, $\delta{=}0.89$ \\
\quad vs.\ advanced classical & 1118 & \textbf{858} & AI, $p{=}.002$ \\
90\,\% interval half-width (cy) & 11509 & \textbf{1957} & 5.9$\times$ tighter \\
Physics-informed hybrid & --- & --- & refuted 2$\times$ \\
\bottomrule
\end{tabular}
\end{center}

The pattern survives contact with reality: on NASA's C-MAPSS benchmark the same pipelines
score 14.9 vs.\ 42.1 cycles RMSE, and the classical side was even upgraded mid-study to a
similarity-based prognostic (1981${\to}$1118) --- AI still wins ($858$, $p{=}0.002$).

\section{The split is physics, not fashion}
AI wins exactly where the task is \emph{statistical aggregation over time}: pulling faint
trends out of noise (detection), integrating a whole degradation history into a forecast
(prognosis), and honestly sizing its own error (uncertainty). It ties exactly where the
\emph{sensor set} cannot distinguish two faults: the cockpit suite's influence-matrix
geometry leaves seven confusable fault pairs at ${\sim}1.3^{\circ}$ separation, and no model
--- classical or learned --- can invert information that never reached the data. Proof:
adding two real station sensors lifts confusable isolation from 0.15 to 0.92
($p{=}0.001$); adding \emph{virtual} (model-derived) versions of the same sensors changes
nothing ($p{=}1.0$). Where learning does add isolation power is fusion: a GRU fusing
multi-flight physics projections more than triples classical multi-point analysis (0.65
vs.\ 0.17) on identical windows, and the ACARS reporting \emph{schedule} itself becomes a
design variable (+79\,\% separability from one extra stabilized-cruise report).

\section{Certificates, not promises}
Two by-products may outlast the scoreboard. An \emph{identifiability certificate} ---
accumulated Fisher information over an engine's actual flight history --- predicts
per-direction diagnosis error \emph{before} deployment, validated against ground truth
(Spearman $\rho{=}0.70$): it tells an operator which faults their fleet can see and which
sensor rescues which blind spot. A \emph{prognostic floor} --- the aleatoric RUL
uncertainty that no model can remove --- shows 87\,\% of early-life error is irreducible,
while at 90\,\% of life the best method still sits $4\times$ above the floor: buy
algorithms for the endgame, calibration for the rest.

\section{What it is worth}
The operational currency is converting unscheduled removals into scheduled ones. Against
the floor-derived ceiling, the AI stack converts a net 35--50\,\% of unscheduled events
across realistic lead-time demands; the traditional stack manages 0--25\,\%. For a
50-aircraft, 100-engine fleet, a Monte-Carlo cost model prices the difference at
\$3.2M/yr median (\$0.9M--\$8.0M p10--p90, ${<}1\,\%$ probability of net loss).

\vspace{1mm}
\noindent\rule{\linewidth}{0.4pt}
{\footnotesize \textbf{Takeaway}: ``AI vs.\ traditional'' is the wrong question. AI is a
strictly better \emph{statistician} on the same data; it is not a new \emph{sensor}. Buy AI
for detection, prognosis, and calibrated uncertainty; buy instrumentation for isolation;
and demand certificates --- identifiability and floor --- that say, in advance, which
regime you are in.\par
\vspace{0.8mm}
Full study: 106-page report, 12 pre-registration tags, every figure and number
auto-generated from versioned artifacts, replicable end-to-end on a desktop
(\texttt{make all}, ${\sim}11$ min). Repository:
\href{https://github.com/cedarcreeks/EHMbrAIn}{github.com/cedarcreeks/EHMbrAIn}.\par}

\end{document}
